% !PRINTED 181
\weberpage{181}

%% [running head, not body text -- commented out 2026-08-07] Page furniture (running
%% head, folio, signature) is excluded from this edition by standing policy.
%% \noindent \S.\,47. \hfill \textit{Normalformen.} \hfill 181

Wir können aber auch umgekehrt schliessen, dass, wenn $S$ einen Normaltheiler vom Index $e$ besitzt, das Formenproblem imprimitiv ist. Wir brauchen nämlich nur, um ein System der Functionen $X_1, \ldots X_m$ zu erhalten, eine Invariante $\varrho$ der Gruppe $T$ zu nehmen, die $e$ verschiedene Werthe in $S$ enthält, und können für $X_1, \ldots X_m$ geradezu die Werthe $\varrho, \varrho_1, \ldots \varrho_{e-1}$ nehmen. Die Gruppe $\Sigma$ ist dann die durch $S$ unter den $\varrho$ hervorgerufene Permutationsgruppe.

Im Sinne des \S.\,44 wird es aber immer darauf ankommen, $m$ so klein als möglich zu machen, und eine Reduction des Problems in diesem Sinne wird nur dann erzielt sein, wenn $m < n$ ist.

\begin{center}
\S.\,47.

\gesperrt{Normalformen.}
\end{center}

Noch eine allgemeine Betrachtung müssen wir anstellen, ehe wir zu speciellen Anwendungen übergehen. Wir haben schon in \S.\,37 gesehen, dass wir aus jeder Gruppe $S$ von linearen Substitutionen unendlich viele isomorphe Gruppen ableiten können, indem wir mit einer willkürlichen Substitution $L$ von derselben Dimension transformiren, also die transformirte Gruppe
\weq{1}{L^{-1} S L}
bilden; und dies ist gleichbedeutend mit der Einführung anderer Variablen $y$ an Stelle von $x$ durch die Substitution
\weq{2}{(x) = L(y).}

Diese Transformation können wir dazu verwenden, um die Gruppe $S$ in einer einfachen Normalform darzustellen, und dann erhalten auch die Invarianten gewisse feste Normalformen, die in manchen Fällen sehr einfache Gestalten annehmen können. Bilden wir für die Normalform die Resolvente des Formenproblems [\S.\,43, (8)], die wir jetzt in der Form schreiben wollen:
\weq{3}{\Phi(\Theta, A_1, A_2, \ldots) = 0,}
worin $A_1, A_2, \ldots$ Invarianten der Gruppe $S$ bedeuten, so wird auch diese, wenn wir die Normalform benutzen, eine einfache Gestalt erhalten. Die Variablen $x_i$ lassen sich, wie wir gesehen
%CONT>>
